\subsection{Definição por Conjuntos}
\label{sec:definicao_conjuntos}

O espaço de trabalho (\emph{workspace}) posicional do manipulador robótico é o conjunto de todas as posições cartesianas que o efetuador pode atingir a partir de configurações válidas - que não colidam com o próprio manipulador robótico - das juntas.  

Esse conjunto pode ser escrito como:
\begin{equation}
\label{eq:workspace_posicional}
W_{pos} = \left\{ \vec{p} \in \mathbb{R}^3 : \exists \, q_j \in Q_{livre} \; \text{com} \; \vec{p} = \vec{p}_E(q) \right\},
\end{equation}

Onde:
\begin{itemize}
    \item $\vec{p}_E(q)$: posição do efetuador obtida pela cinemática direta para uma dada configuração de juntas $q$;
    \item $Q_{livre}$: conjunto de juntas admissíveis sem colisões ou violações de restrições, como o Keep-Out Zones;
    \item $\vec{p} \in \mathbb{R}^3$: vetor posição cartesiano atingível pelo efetuador.
\end{itemize}

O conjunto de juntas livres de colisão é dado por:
\begin{equation}
\label{eq:workspace_Qlivre}
Q_{livre} = Q_{adm} - Q_{col}.
\end{equation}

Onde:
\begin{itemize}
    \item $Q_{adm}$: conjunto admissível de juntas, definido pelos limites geométricos, cinemáticos e dinâmicos apresentados na seção \eqref{sec:conjunto_juntas};
    \item $Q_{col}$: conjunto de configurações inválidas devido a auto-colisão do manipulador robótico ou violação das restrições de segurança, por exemplo, Keep-Out Zones.
\end{itemize}

Portanto, o workspace posicional $W_{pos}$ corresponde à região do espaço cartesiano tridimensional que pode ser alcançada pelo efetuador respeitando todas as restrições operacionais do sistema.
